%-----------------------------------------------------------------------
%  02 -- the modified files, grouped by the role they play. Six frames.
%
%  Each file is shown as ours -> the checkpoint69m file it shadows, with the
%  path linked to GitHub, and answers three questions; the fourth (why it
%  still fits MITgcm's structure) is shared across the group and answered by
%  the closing strip.
%
%  Difference sizes are `diff <c69m file> <our file>` changed-line counts
%  measured against the vendored tree, not quoted from the change note.
%  \input by main.tex. Not compiled alone.
%-----------------------------------------------------------------------

\begin{frame}{The modified files 1/6 --- the two call sites}
  \framesubtitle{\code{model/src/} --- where the hooks are invoked}

\Gap
Tapenade generates a reverse-sweep call only where active data crosses an
interface. Both hooks are called with three passive scalars, so both are
absent from the generated backward sweep.

\pairfile{code_tap/forward_step.F}{../../MITgcm/model/src/forward_step.F}{25 lines}
{The forward time-stepping driver. It calls \code{DUMMY_IN_STEPPING} in its
\code{ALLOW_AUTODIFF_MONITOR} block, and the two \code{AUTODIFF_INADMODE_*}
hooks at the start and end of each step.}
{Nothing in those calls is differentiable, so no reverse-sweep counterpart is
produced for any of them.}
{Under \cpp{ifdef ALLOW_TAPENADE} the same three calls also pass model state:
eleven fields to the dump hook, \code{uVel} to each mode switch. The
three-argument calls remain in \cpp{else}.}

\pairfile{code_tap/integr_continuity.F}{../../MITgcm/model/src/integr_continuity.F}{8 lines}
{Integrates the continuity equation and updates the free surface. It calls
\code{DUMMY_FOR_ETAN} from a block of its own because the free-surface adjoint
is half a time step out of phase with the fields in \code{forward_step.F}.}
{The generated \code{integr_continuity_b.f} retained an \code{EXTERNAL}
declaration and no call, so no \code{ADJetan} field was ever written.}
{One guarded call site passing \code{etaN}.}

\Safe{Four guarded call sites and no other change. The \cpp{else} branch is
the checkpoint69m line, so a TAF or OpenAD build preprocesses to the unmodified
source.}
\end{frame}


\begin{frame}{The modified files 2/6 --- the two dump-hook interfaces}
  \framesubtitle{\code{pkg/autodiff/} --- the no-op routines the call sites reach}

\Gap
Both routines are deliberately empty and declare no fields, so there is no
interface through which activity can pass. Under TAF this does not matter: a
\code{.flow} directive attaches a hand-written adjoint to the hook regardless.

\pairfile{code_tap/dummy_in_stepping.F}{../../MITgcm/pkg/autodiff/dummy_in_stepping.F}{31 lines}
{An empty subroutine on \code{(myTime, myIter, myThid)}, so the call in
\code{forward_step.F} resolves in forward runs and checkpoint re-forwards;
\code{ADNAME} binds \code{ADDUMMY_IN_STEPPING} to it for TAF.}
{No field arguments, so nothing can be declared active to Tapenade.}
{The \code{SUBROUTINE} statement and declarations carry the eleven dump fields
under the \code{ALLOW_TAPENADE} guard, with the checkpoint69m interface in
\cpp{else}. The body remains empty.}

\pairfile{code_tap/dummy_for_etan.F}{../../MITgcm/pkg/autodiff/dummy_for_etan.F}{20 lines}
{The same, for the free surface.}
{As above.}
{One field, \code{etaN}, added under the same guard.}

\Safe{This is the interface-widening idiom \code{pkg/autodiff} already applies
under \code{AUTODIFF_TAMC_COMPATIBILITY}. Neither routine appears in any
package's \code{*_ad_diff.list}, so no AD tool differentiates them and both
remain no-ops.}
\end{frame}


\begin{frame}{The modified files 3/6 --- the adjoint-mode switches}
  \framesubtitle{\code{pkg/autodiff/} --- entering and leaving adjoint mode}

\Gap
MITgcm can run the reverse sweep with different physics from the forward
trajectory (\code{inAdMode}, \code{\vfin}) --- the usual way of keeping a long
adjoint from diverging. The routines applying those settings are TAF-named, so
nothing Tapenade generated called them: \code{\vfin{} = 10} could be set in the
namelist with no code path applying it.

\pairfile{code_tap/autodiff_inadmode_set.F}{../../MITgcm/pkg/autodiff/autodiff_inadmode_set.F}{20 lines}
{An empty hook marking where the backward sweep enters adjoint mode.}
{No active argument, so no \code{_B} routine is generated and the settings are
never applied.}
{\code{uVel} added under the \code{ALLOW_TAPENADE} guard purely to carry
activity across the interface; it is never read.}

\pairfile{code_tap/autodiff_inadmode_unset.F}{../../MITgcm/pkg/autodiff/autodiff_inadmode_unset.F}{20 lines}
{The matching hook for leaving adjoint mode.}
{As above.}
{The same argument, so the restore runs at each backward step's end and
checkpoint re-forwards use forward physics.}

\Safe{Confirmed an exact no-op when adjoint and forward settings agree
(\code{\vfin{} = viscFacInFw = 1}): run 31024 reproduces 31023 bitwise.}
\end{frame}


\begin{frame}{The modified files 4/6 --- bodies for routines shipped as stubs}
  \framesubtitle{\code{pkg/tapenade/} --- where checkpoint69m provides names but no implementation}

\pairfile{code_tap/dummy_tap.F}{../../MITgcm/pkg/tapenade/dummy_tap.F}{456 lines}
{Provides \code{DUMMY_IN_STEPPING_B}, \code{DUMMY_FOR_ETAN_B} and their
tangent-mode counterparts as empty routines on the three passive scalars, so
that a Tapenade build links.}
{The implementations. Even if a reverse-sweep call were generated, it would
reach a routine that returns immediately --- which is why the Tapenade path
produces no \code{ADJ*} output as distributed.}
{Full bodies at the wider signatures: the adjoint arguments are folded across
tile boundaries and written through \code{DUMP_ADJ_XYZ}/\code{DUMP_ADJ_XY}, with
\code{ADJetan} on its own record counter. Two short wrappers call MITgcm's
existing \code{ADAUTODIFF_INADMODE_SET}/\code{UNSET}, so the parameter logic is
not duplicated.}

\pairfile{code_tap/stubs_tap_adj.F}{../../MITgcm/pkg/tapenade/stubs_tap_adj.F}{183 lines}
{Provides the five \code{ADEXCH_*} adjoint halo exchanges as routines that
print \code{"Called not yet defined"} and return.}
{The implementations. The dump path calls them to fold tile-halo adjoint
contributions back into the interior cells that own them before writing.}
{Implemented using the same \code{EXCH2_*_CUBE_AD} routines the adjoint
dynamics already calls. The effect is confined to the dumps: \code{fc} and the
\code{adxx_*} gradients were correct throughout.}

\Safe{Both files complete routines checkpoint69m already declares. No
interface changes, no file is displaced, and 639 of the contribution's lines
replace \code{RETURN} statements.}
\end{frame}


\begin{frame}{The modified files 5/6 --- the Tapenade configuration}
  \framesubtitle{\code{tools/} --- the declaration the whole mechanism rests on}

\pairfile{code_tap/flow_tap_local}{../../MITgcm/tools/TAP_support/flow_tap}{4 stanzas}
{Tapenade's external library. It declares each of the four hooks with three
parameters, every one \code{ReadNotWritten} --- that is, entirely passive.}
{This declaration is precisely what tells Tapenade the hooks carry no
derivative information, and so is the reason they are removed from the reverse
sweep. It is MITgcm's own description of the routines, and it is accurate for
the interfaces as they stand.}
{Four stanzas re-declaring the same routine names with their field arguments
\code{ReadThenWritten}: fourteen parameters for the dump hook, four for the
others. This one declaration is what causes the \code{_B} calls to be
generated; the source changes are the interface work around it.}

\pairfile{code_tap/adjoint_tap_local}{../../MITgcm/tools/adjoint_options/adjoint_tap}{2 lines}
{The \code{-adof} options file \code{genmake2} sources to assemble
\code{TAPENADE_FLAGS}, including \code{-ext tools/TAP_support/flow_tap}.}
{Nothing structural. This file exists only so that the vendored MITgcm tree
need not be edited in place.}
{Sources the stock options file and appends the second \code{-ext} library.
The order matters: Tapenade retains the \emph{last} declaration of an external,
and when the local stanzas were supplied first they were overridden by the
stock ones (\code{TC32 Conflicting numbers of arguments}), with no \code{_B}
call generated.}

\Safe{\code{flow_tap} is read by Tapenade alone, so neither TAF nor OpenAD is
affected. In a contribution the stanzas would be edited into \code{flow_tap}
and the wrapper would not be needed; \code{genmake2} is unmodified either way.}
\end{frame}


\begin{frame}{The modified files 6/6 --- the adjoint-mode parameter bodies}
  \framesubtitle{\code{pkg/autodiff/} --- supplied for our experiments, and not part of the proposal}

\pairfile{code_tap/variants/adjointViscosity/autodiff_inadmode_set_ad.F}{../../MITgcm/pkg/autodiff/autodiff_inadmode_set_ad.F}{applies}
{Holds \code{ADAUTODIFF_INADMODE_SET}, which applies the adjoint-mode
parameters when the backward sweep begins.}
{Reachability, addressed on frame 7 --- and a body raising viscosity and
diffusivity for the reverse sweep.}
{The checkpoint69m body together with the block that applies the elevated
values, supplied as a second \code{-mods} directory and compiled only by the
corresponding build.}

\pairfile{code_tap/variants/adjointViscosity/autodiff_inadmode_unset_ad.F}{../../MITgcm/pkg/autodiff/autodiff_inadmode_unset_ad.F}{restores}
{Holds the matching routine for leaving adjoint mode.}
{The restore side had no implementation anywhere: the parameters are declared
and read, and nothing applied them. Because the hook was unreachable, this had
gone unnoticed.}
{A restore block that returns the forward values at each backward step's end,
so that checkpoint re-forwards integrate with forward physics.}

\Strip{ImpBlue}{Scope}{These two bodies encode the viscosity settings of one
set of experiments, so they are configuration rather than mechanism and are not
part of what we would propose. The general half --- the hooks that make them
reachable at all --- is on frame 7.}
\end{frame}
